\documentclass[letterpaper,11pt]{article}
\usepackage[empty]{fullpage}
\usepackage[hidelinks]{hyperref}
\usepackage[english]{babel}
\addtolength{\oddsidemargin}{-0.6in}
\addtolength{\evensidemargin}{-0.5in}
\addtolength{\textwidth}{1.19in}
\addtolength{\topmargin}{-.9in}
\addtolength{\textheight}{1.7in}
\urlstyle{same}
\raggedbottom
\raggedright
\setlength{\parindent}{0pt}
\setlength{\parskip}{10pt}
\pagestyle{empty}

\begin{document}

\begin{center}
    {\LARGE \scshape Jendrick Rian Galut} \\ \vspace{2pt}
    \small Surrey, BC ~$|$~
    672-272-3193 ~$|$~
    \href{mailto:rcgalut1@gmail.com}{\underline{rcgalut1@gmail.com}} ~$|$~
    \href{https://www.linkedin.com/in/jendrick-rian-galut-374310307/}{\underline{linkedin.com/in/jendrick-galut}} ~$|$~
    \href{https://github.com/Jeam-cloud}{\underline{github.com/jeam-cloud}}
\end{center}

\vspace{8pt}

August 12, 2026

Hiring Team \\
Sanctuary \\
Vancouver, BC

Dear Hiring Team,

I am writing to apply for the Technical Product Management Internship starting September 2026. I am a Computing Science student at Trinity Western University, and what draws me to this role is not robotics in the abstract but the specific discipline the posting describes: taking a technical system, pinning down what it actually needs to do, writing that down so it is traceable, and checking afterward that it still does it.

That is the part of building software I have gravitated toward on my own. On ARTHUR, a local-first AI assistant I am currently building, the hardest problem was not writing code but defining what the system was and was not allowed to do: I translated a set of safety requirements into a concrete design (task-mode privilege separation, an approval broker that gates risky actions, a sandboxed execution environment with per-tool network policy), and when I identified a gap -- the system was exposed to prompt injection from untrusted content -- I specified the fix, implemented it, and validated with a dedicated test suite that the requirement actually held. On RateMyTWU, a professor-review platform I scoped and shipped that is now used organically by 50+ students, I made the feature-set decisions myself: what to build first, what a degree-planning feature needed to do, and how to sequence it against a small team of one.

I have also spent the last several months doing requirements work in a much less technical setting. As a Building Manager, I track every incident, maintenance request, and bylaw issue through a structured case-management system, and the job only works if that documentation is traceable enough that the next person handling it -- a contractor, a board member, a resident -- doesn't have to ask me what happened. As a Teaching Assistant across two computing science courses, I have spent two semesters translating technical material for people who did not build it, which is close to what I understand this role to require when connecting engineering trade-offs to customer and business needs.

One thing I want to raise directly rather than leave for you to notice: my expected graduation date is April 2028, later than the "final year or recent graduate" description in the posting. I am including it because I would rather you have accurate information than a polished application that glosses over it. If there is flexibility here, or if a different term would be a better fit, I would welcome that conversation.

Thank you for considering my application. I would be glad to discuss how my background fits the role.

Sincerely,

Jendrick Rian Galut

\end{document}
